\documentclass[exercice]{thedocument}%
\startdatakeys
\source{}
\BO{2026}
\cycle{3}
\competencesDuSocle{RE}
\objectifsApprentissage{C2ai,C2aj,C2ak,C2al,C2am}
\automatismes{B3ab,B3aa}
\enddatakeys
\begin{document}%
À l'aide d'un théodolite (appareil servant à mesurer des angles) et d'un
décamètre, un géomètre a fait les relevés suivants.
\begingroup\centering\pdf{6eme_Indigo_214_44_picture}\par\endgroup
Il voudrait connaître la distance entre la grange et le puit, qu'il ne peut pas
mesurer directement à cause de la rivière.
\vskip10pt
Tracer cette figure (1 cm sur le dessin représentant 1 m en réalité) et
donner une estimation de la distance entre la grange et le puits.
\end{document}
